\section{受控连接与“编织”之严格化：控制输入、可达性与优化问题}

\subsection{控制输入与受控转移核}

在第 9 节离散连接语义下，平行移动 $\mathcal P_{x\to x'}$ 与允许边集 $\mathcal E_m$ 共同决定提升路径族。将“编织/选择”形式化为控制输入序列 $\{u_t\}$，其作用可取两类等价形式：
\begin{enumerate}
  \item \textbf{受控允许边：} $\mathcal E_m(u_t)\subseteq X_m\times X_m$；
  \item \textbf{受控转移核：} $K_t^{u_t}(z\to z')$。
\end{enumerate}
于是演化写为
$$
z_{t+1}\sim K_t^{u_t}(z_t,\cdot),\qquad z_t=(x_t,r_t)\in E_m.
$$
控制输入并不破坏局部性公理：允许 $K_t^{u_t}$ 仅依赖有限邻域信息与有限记忆量。

\subsection{控制目标与代价泛函}

给定目标一致性 $\mathcal G$ 或终端约束 $\mathcal F_T$，可定义代价泛函
$$
J(\pi)=\mathbb E^{\pi}\left[\sum_{t=0}^{T-1} c(z_t,u_t)+\lambda\,\Phi_T(z_T)\right],
$$
其中 $\pi$ 为策略（从局部可观测信息到控制的映射），$c$ 可包含：
\begin{itemize}
  \item 残差更新成本（与迭代复杂度相关）；
  \item 接口缺陷惩罚（与第 7 节耦合项相关）；
  \item holonomy 复杂度代理量（与第 9 节曲率相关）；
  \item 分辨率调度成本（与第 11 节 AMR 相关）。
\end{itemize}
该形式使“愿力/选择”的叙述可被替换为可检验的控制论问题：在局部约束与有限资源下寻找最小化 $J(\pi)$ 的策略族。

\subsection{双向搜索的算法化表述}

在仅前向推进的情形，寻找满足终端约束的轨道等价于在扩展态图上进行可达性搜索，其复杂度由状态空间规模与回路结构控制。引入后向权重 $C_t$（第 13 节）等价于对同一图进行前向--后向消息传递，其在树状或近树状结构上可显著降低搜索分支，且在一般图上提供可行的变分/截断近似。该结论与 HMM 平滑在序列模型中的复杂度优势一致。\cite{rabiner_hmm_tutorial}

\subsection{连接可控性与 holonomy 类选择}

在固定 $X_m$ 与 $R_m$ 的条件下，控制输入可被解释为在允许边集与平行移动族之间选择，使得某些闭路的 holonomy 被压制或增强，从而改变历史依赖结构对时间成本 $T(\mathcal G)$ 的影响。该观点将“编织时间纤维”转写为“在离散连接上选择提升路径族，使 holonomy 谱的有效复杂度下降”。
